% Terjemahan Bahasa Indonesia dari chapter2.tex baris 721--910.
% Sumber: Methods of Algebra, Volume 2, commit
% 9a5803ff77dd3257484cb177f851a73770a59dd3 (CC BY 4.0).
% stable-unit-id: o014.aljabr2.chapter2.direct-sum-decomposition

% segment-id: o014.li2.chapter2.direct-sum-decomposition.h001
\section{Dekomposisi Jumlah Langsung}\label{sec:direct-sum}
\hypertarget{o014.aljabr2.chapter2.direct-sum-decomposition}{ }

% segment-id: o014.li2.chapter2.direct-sum-decomposition.p001
Pertama-tama, kita membahas cara merepresentasikan morfisme antarjumlah
langsung dengan matriks. Argumennya persis sama seperti dalam
teori modul \cite[\S 6.3]{Li1}; di sini kita hanya memberikan tinjauan singkat.

% segment-id: o014.li2.chapter2.direct-sum-decomposition.p002
Tetapkan suatu kategori aditif $\mathcal{A}$. Tinjau objek-objek
$X_1, \ldots, X_n$ dan $X'_1, \ldots, X'_m$ dari $\mathcal{A}$. Jumlah
langsung $\bigoplus_{i=1}^n X_i$ yang diperkenalkan dalam
Definisi~\ref{def:additive-category} dilengkapi dengan keluarga morfisme
$X_j \xrightarrow{\iota_j} \bigoplus_{i=1}^n X_i \xrightarrow{p_j} X_j$.
Demikian pula, untuk $\bigoplus_{i=1}^m X'_i$ terdapat $p'_j,\iota'_j$, dan
seterusnya. Sifat universal produk dan koproduk memberikan
% segment-id: o014.li2.chapter2.direct-sum-decomposition.q001
\begin{align*}
	\Hom \left( \bigoplus_{j=1}^n X_j, \bigoplus_{i=1}^m X'_i \right) & \xrightarrow[\sim]{\phi \mapsto (p'_i \phi)_i } \prod_{i=1}^m \Hom\left(\bigoplus_{j=1}^n X_j , X'_i \right) \\
	& \xrightarrow[\sim]{(p'_i \phi)_i \mapsto (p'_i \phi \iota_j)_{i, j}} \prod_{j=1}^n \prod_{i=1}^m  \Hom(X_j, X'_i).
\end{align*}
Dengan demikian, setiap morfisme
$\phi: \bigoplus_{j=1}^n X_j \to \bigoplus_{i=1}^m X'_i$ bersesuaian dengan
matriks
% segment-id: o014.li2.chapter2.direct-sum-decomposition.q002
\[ \mathcal{M}(\phi) = (\phi_{ij})_{\substack{1 \leq i \leq m \\ 1 \leq j \leq n}} = \begin{pmatrix}
	\phi_{11} & \cdots & \phi_{1n} \\
	\vdots & \ddots & \vdots \\
	\phi_{m1} & \cdots & \phi_{mn}
\end{pmatrix}, \quad \phi_{ij} := p'_i \phi \iota_j: X_j \to X'_i. \]

% segment-id: o014.li2.chapter2.direct-sum-decomposition.p003
Sebaliknya, setiap matriks $\mathcal{M} = (\phi_{ij})_{\substack{ 1 \leq i \leq m \\ 1 \leq j \leq n }}$
yang tersusun dari $\phi_{ij}: X_j \to X'_i$ secara unik menentukan morfisme
$\phi$ sedemikian sehingga $\mathcal{M}=\mathcal{M}(\phi)$. Komposisi morfisme
bersesuaian dengan perkalian matriks:
% segment-id: o014.li2.chapter2.direct-sum-decomposition.q003
\[ \mathcal{M}(\psi\phi) = \mathcal{M}(\psi) \mathcal{M}(\phi) \]
dengan perkalian entri-entri matriks diberikan oleh komposisi
$\Hom(X'_j, X''_i) \times \Hom(X_k, X'_j) \to \Hom(X_k, X''_i)$.

% segment-id: o014.li2.chapter2.direct-sum-decomposition.le001
\begin{lema}\label{prop:isomorphism-matrix}
	Diberikan keluarga morfisme $\phi_i:X_i\to X'_i$ dalam $\mathcal{A}$,
	dengan $i=1,\ldots,n$, tetapkan
	$\phi := (\phi_1, \ldots, \phi_n): \bigoplus_{i=1}^n X_i \to
	\bigoplus_{i=1}^n X'_i$. Maka $\phi$ merupakan isomorfisme jika dan hanya
	jika setiap $\phi_i$ merupakan isomorfisme.
\end{lema}
% segment-id: o014.li2.chapter2.direct-sum-decomposition.pf001
\begin{bukti}
	Implikasi $\impliedby$ sudah jelas. Untuk $\implies$, tuliskan matriks
	$\mathcal{M}(\phi^{-1})$ sebagai $(\psi_{ij})_{i,j}$. Maka
	% segment-id: o014.li2.chapter2.direct-sum-decomposition.q004
	\[ \mathcal{M}(\phi^{-1}) \mathcal{M}(\phi) = \begin{pmatrix}
		\psi_{11} \phi_1 & \cdots & \psi_{1n} \phi_n \\
		\vdots & \ddots & \vdots \\
		\psi_{n1} \phi_1 & \cdots & \psi_{nn} \phi_n
	\end{pmatrix} = \begin{pmatrix}
		\identity_{X_1} & & \\
		& \ddots & \\
		& & \identity_{X_n}
	\end{pmatrix}\]
	sehingga $\phi_1,\ldots,\phi_n$ semuanya memiliki invers kiri. Dengan cara
	yang sama, semuanya memiliki invers kanan.
\end{bukti}

% segment-id: o014.li2.chapter2.direct-sum-decomposition.p004
Masih dengan dekomposisi jumlah langsung
$X \simeq \bigoplus_{i=1}^n X_i$, dengan $n\geq 1$ dan
$X_i\neq 0$ untuk setiap $i$; darinya diperoleh keluarga morfisme
$\iota_i:X_i\to X$ dan $p_i:X\to X_i$. Untuk setiap $1\leq i\leq n$,
tetapkan $e_i:=\iota_i p_i\in\End(X)$. Morfisme-morfisme ini mempunyai
sifat-sifat berikut. Ingat bahwa suatu \emph{idempoten} dalam gelanggang
$\End(X)$ berarti unsur $e\in\End(X)$ yang memenuhi $e^2=e$; jika tidak
menimbulkan kerancuan, kita juga menyebut $e$ sebagai idempoten dalam
$\mathcal{A}$.
\index{idempoten (idempotent)}

% segment-id: o014.li2.chapter2.direct-sum-decomposition.l001
\begin{enumerate}[(E1)]
	% segment-id: o014.li2.chapter2.direct-sum-decomposition.i001
	\item Setiap $e_i$ merupakan idempoten:
	$e_i^2=(\iota_i p_i)(\iota_i p_i)=\iota_i(p_i\iota_i)p_i=\iota_i p_i$.
	% segment-id: o014.li2.chapter2.direct-sum-decomposition.i002
	\item Ortogonalitas: $i\neq j \implies
	e_i e_j=\iota_i p_i\iota_j p_j=0$.
	% segment-id: o014.li2.chapter2.direct-sum-decomposition.i003
	\item Persamaan $\sum_{i=1}^n e_i=1$ berlaku dalam gelanggang $\End(X)$.
\end{enumerate}

% segment-id: o014.li2.chapter2.direct-sum-decomposition.p005
Mulai sekarang, kita selalu memandang suku-suku jumlah langsung $X_i$ sebagai
subobjek dari $X$ dan menuliskan dekomposisinya sebagai persamaan
$X=\bigoplus_{i=1}^n X_i$. Sebaliknya, jika $\mathcal{A}$ merupakan kategori
abelian, suatu keluarga idempoten
$e_1,\ldots,e_n\in\End(X)$ yang memenuhi (E1)---(E3) menentukan subobjek-subobjek
dari $X$
% segment-id: o014.li2.chapter2.direct-sum-decomposition.q005
\[ X_i := \Ker\left( \sum_{j \neq i} e_j \right) = \Image(e_i), \quad i=1, \ldots, n. \]
Di satu pihak, subobjek-subobjek tersebut dilengkapi dengan monomorfisme
$X_i\xrightarrow{\iota_i}X$; di pihak lain, $e_i:X\to X$ secara unik terurai
sebagai $X\xrightarrow{p_i}X_i\xrightarrow{\iota_i}X$. Inilah tepatnya data
yang diperlukan untuk jumlah langsung.

% segment-id: o014.li2.chapter2.direct-sum-decomposition.p006
Peralihan dari idempoten-idempoten $e_1,\ldots,e_n$ ke suku-suku jumlah
langsung $X_1,\ldots,X_n$ tidak memerlukan seluruh sifat kategori abelian;
kita hanya perlu mengasumsikan bahwa $\mathcal{A}$ merupakan
kategori-$\cate{Ab}$ dengan objek nol dan bahwa setiap idempoten mempunyai
kernel. Kategori semacam itu disebut \emph{kategori Karoubi} atau
\emph{kategori pseudoabelian}. Hal ini dibahas lebih lanjut dalam latihan bab
ini.
\index{kategori!Karoubi (Karoubian)}

% segment-id: o014.li2.chapter2.direct-sum-decomposition.pr001
\begin{proposisi}\label{prop:idempotent-decomposition}
	Diberikan objek $X\neq 0$ dari kategori abelian (atau, secara lebih umum,
	kategori Karoubi) $\mathcal{A}$ dan $n\in\Z_{\geq 1}$, konstruksi di atas
	memberikan bijeksi
	% segment-id: o014.li2.chapter2.direct-sum-decomposition.g001
	\[\begin{tikzcd}[row sep=small]
		\left\{\begin{gathered}
			\text{dekomposisi jumlah langsung}\\[-2pt]
			X = \bigoplus_{i=1}^n X_i
		\end{gathered}\right\}
		\arrow[leftrightarrow, r, "1:1"] &
		\left\{\begin{gathered}
			(e_i)_{i=1}^n \in \End(X)^n:\\[-2pt]
			\text{memenuhi (E1)---(E3)}
		\end{gathered}\right\}.
	\end{tikzcd}\]
	Pengurutan ulang suku-suku jumlah langsung $X_1,\ldots,X_n$ bersesuaian
	dengan pengurutan ulang $e_1,\ldots,e_n$.
\end{proposisi}
% segment-id: o014.li2.chapter2.direct-sum-decomposition.pf002
\begin{bukti}
	Lihat \cite[Proposisi 6.12.4]{Li1}.
\end{bukti}

% segment-id: o014.li2.chapter2.direct-sum-decomposition.p007
Dalam kasus khusus $n=2$, dekomposisi jumlah langsung berkaitan dengan suatu
kelas khusus barisan eksak pendek, yang disebut barisan eksak pendek terbelah.

% segment-id: o014.li2.chapter2.direct-sum-decomposition.pr002
\begin{proposisi}\label{prop:split-ses}
	\index{barisan eksak pendek!terbelah (split)}
	Untuk barisan eksak pendek
	$0\to X'\xrightarrow{f}X\xrightarrow{g}X''\to 0$ dalam suatu kategori
	abelian, pernyataan-pernyataan berikut ekuivalen.
	% segment-id: o014.li2.chapter2.direct-sum-decomposition.l002
	\begin{enumerate}[(i)]
		% segment-id: o014.li2.chapter2.direct-sum-decomposition.i004
		\item Terdapat $s:X''\to X$ sedemikian sehingga
		$gs=\identity_{X''}$.
		% segment-id: o014.li2.chapter2.direct-sum-decomposition.i005
		\item Terdapat $r:X\to X'$ sedemikian sehingga
		$rf=\identity_{X'}$.
		% segment-id: o014.li2.chapter2.direct-sum-decomposition.i006
		\item Terdapat diagram
		% segment-id: o014.li2.chapter2.direct-sum-decomposition.g002
		$\begin{tikzcd}
			X' \arrow[yshift=-0.7ex, r, "f"'] & X \arrow[yshift=0.7ex, l, "r"'] \arrow[yshift=-0.7ex, r, "g"'] & X'' \arrow[yshift=0.7ex, l, "s"']
		\end{tikzcd}$
		sedemikian sehingga $X\simeq X'\oplus X''$; lihat tinjauan tentang
		biproduk dalam \S\ref{sec:Ker-Coker}.
		% segment-id: o014.li2.chapter2.direct-sum-decomposition.i007
		\item Pemetaan $g_*:\Hom(S,X)\to\Hom(S,X'')$ (yang memetakan
		$\varphi$ ke $g\varphi$) bersifat surjektif untuk setiap objek $S$.
		% segment-id: o014.li2.chapter2.direct-sum-decomposition.i008
		\item Pemetaan $f^*:\Hom(X,S)\to\Hom(X',S)$ (yang memetakan $\psi$
		ke $\psi f$) bersifat surjektif untuk setiap objek $S$.
	\end{enumerate}
	Jika salah satu kondisi di atas berlaku, barisan eksak pendek
	$0\to X'\to X\to X''\to 0$ disebut \emph{terbelah}. Morfisme
	$s:X''\to X$ dalam (i), $r:X\to X'$ dalam (ii),\footnote{Catatan
	penerjemah (O014-C024): sumber menulis $r:X'\to X$ di sini, tetapi (ii),
	persamaan $rf=\identity_{X'}$, dan konstruksi berikutnya semuanya
	menentukan $r:X\to X'$. Target menggunakan arah yang bertipe benar.} dan
	isomorfisme $X\rightiso X'\oplus X''$ dalam (iii) (yang dilambangkan dengan
	$\Phi$) saling bersesuaian sebagai berikut:
	% segment-id: o014.li2.chapter2.direct-sum-decomposition.l003
	\begin{compactitem}
		% segment-id: o014.li2.chapter2.direct-sum-decomposition.i009
		\item dari $\Phi$ ke $s$: ambil komposisi
		$X''\xrightarrow{\iota_2}X'\oplus X''\xrightarrow{\Phi^{-1}}X$;
		% segment-id: o014.li2.chapter2.direct-sum-decomposition.i010
		\item dari $s$ ke $r$: morfisme $\identity_X-sg:X\to X$ secara unik
		terurai sebagai $X\xrightarrow{r}X'\xrightarrow{f}X$; inilah $r$ yang
		dikehendaki;
		% segment-id: o014.li2.chapter2.direct-sum-decomposition.i011
		\item dari $r$ ke $\Phi$: ambil $(r,g):X\to X'\oplus X''$, yang
		merupakan suatu isomorfisme.
	\end{compactitem}

	Selain itu, untuk $s$ yang diberikan atau $r$ yang bersesuaian dengannya,
	idempoten-idempoten ortogonal yang bersesuaian dengan dekomposisi jumlah
	langsung dalam (iii) adalah
	% segment-id: o014.li2.chapter2.direct-sum-decomposition.q006
	\[ e' := fr = \identity_X - sg, \quad e'' := sg = \identity_X - fr. \]
\end{proposisi}
% segment-id: o014.li2.chapter2.direct-sum-decomposition.pf003
\begin{bukti}
	Untuk kasus modul, lihat \cite[Proposisi 6.8.5]{Li1}. Karena buktinya hanya
	menggunakan operasi aljabar dalam himpunan-himpunan $\Hom$ dan karakterisasi
	biproduk, yang berlaku pula dalam kategori abelian, argumennya dapat
	diterapkan tanpa perubahan. Rinciannya tidak diulangi di sini.
\end{bukti}

% segment-id: o014.li2.chapter2.direct-sum-decomposition.p008
Dalam situasi Proposisi~\ref{prop:split-ses}, diberikan $f:X'\to X$ dan
$g:X\to X''$, morfisme $s$ disebut \emph{penampang} dari $g$ jika
$gs=\identity_{X''}$, sedangkan $r$ disebut \emph{retraksi} dari $f$
jika $rf=\identity_{X'}$; istilah-istilah ini berasal dari topologi.
Keberadaan penampang dan retraksi masing-masing mengakibatkan bahwa $g$
merupakan epimorfisme dan $f$ merupakan monomorfisme, tetapi kebalikannya tidak
berlaku.
\index{penampang (section)}
\index{retraksi (retract)}

% segment-id: o014.li2.chapter2.direct-sum-decomposition.p009
Perhatikan pula bahwa kelima pernyataan (i)---(v) dalam proposisi tersebut
bersifat swa-dual.

% segment-id: o014.li2.chapter2.direct-sum-decomposition.d001
\begin{definisi}\label{def:indecomposable}\index{objek!tak terurai (indecomposable)}
	Misalkan $S$ merupakan objek tak nol dalam kategori abelian $\mathcal{A}$.
	Jika $S=S'\oplus S''$ mengakibatkan $S'=0$ atau $S''=0$, maka $S$ disebut
	\emph{objek tak terurai}.
\end{definisi}

% segment-id: o014.li2.chapter2.direct-sum-decomposition.co001
\begin{korolari}\label{prop:indecomposable-criterion}
	Objek tak nol $X$ dalam kategori abelian merupakan objek tak terurai jika
	dan hanya jika
	% segment-id: o014.li2.chapter2.direct-sum-decomposition.q007
	\[ \forall e \in \End(X), \quad e^2 = e \iff (e = 0 \;\vee\; e = 1). \]
\end{korolari}
% segment-id: o014.li2.chapter2.direct-sum-decomposition.pf004
\begin{bukti}
	Terapkan Proposisi~\ref{prop:idempotent-decomposition}.
\end{bukti}

% segment-id: o014.li2.chapter2.direct-sum-decomposition.p010
Kita ingin mempelajari dekomposisi berbentuk
$X=X_1\oplus\cdots\oplus X_n$, dengan $X\neq 0$ dan setiap $X_i$ merupakan
objek tak terurai. Ada dua persoalan: keberadaan dan keunikan. Untuk kasus
modul, inilah isi Teorema Krull--Remak--Schmidt; lihat
\cite[Korolari 6.12.9]{Li1}. Untuk kategori abelian umum, kita memerlukan
sejumlah persiapan. Pembahasan berikut mengikuti pendekatan \cite{Kr15}.

% segment-id: o014.li2.chapter2.direct-sum-decomposition.d002
\begin{definisi}\label{def:bichain}
	\index{syarat rantai ganda (bi-chain condition)}
	Misalkan $\mathcal{A}$ merupakan sebarang kategori. Suatu \emph{rantai
	ganda} di dalamnya adalah data $(X_n,\alpha_n,\beta_n)_{n=0}^\infty$,
	dengan $X_n$ objek dari $\mathcal{A}$ dan
	% segment-id: o014.li2.chapter2.direct-sum-decomposition.g003
	$\begin{tikzcd}
		X_n \arrow[twoheadrightarrow, r, shift left, "\alpha_n"] & X_{n+1} \arrow[hookrightarrow, l, shift left, "\beta_n"]
	\end{tikzcd}$
	merupakan pasangan morfisme di antara objek-objek tersebut, dengan
	$\alpha_n$ epimorfisme dan $\beta_n$
	monomorfisme ($n\in\Z_{\geq 0}$). Suatu objek $X$ dari $\mathcal{A}$
	dikatakan memenuhi \emph{syarat rantai ganda} jika, untuk setiap rantai
	ganda $(X_n,\alpha_n,\beta_n)_{n=0}^\infty$ dengan $X_0=X$, morfisme
	$\alpha_n$ dan $\beta_n$ merupakan isomorfisme untuk $n\gg 0$.
\end{definisi}

% segment-id: o014.li2.chapter2.direct-sum-decomposition.p011
Hasil berikut merupakan perumuman dari \cite[Lema 6.11.5]{Li1}.
% segment-id: o014.li2.chapter2.direct-sum-decomposition.le002
\begin{lema}[Lema Fitting dalam Kategori Abelian]\label{prop:Abel-cat-Fitting}
	Misalkan $\mathcal{A}$ merupakan kategori abelian, $X$ merupakan objek tak
	nol di dalamnya yang memenuhi syarat rantai ganda, dan $f\in\End(X)$.
	% segment-id: o014.li2.chapter2.direct-sum-decomposition.l004
	\begin{enumerate}[(i)]
		% segment-id: o014.li2.chapter2.direct-sum-decomposition.i012
		\item Untuk $n\gg 0$, $\Ker(f^n)$ dan $\Image(f^n)$ tidak bergantung
		pada $n$; keduanya masing-masing dilambangkan dengan $\Ker(f^\infty)$
		dan $\Image(f^\infty)$. Kita mempunyai
		$X=\Ker(f^\infty)\oplus\Image(f^\infty)$.
		% segment-id: o014.li2.chapter2.direct-sum-decomposition.i013
		\item Jika $X$ tak terurai, maka $f$ bersifat invertibel atau nilpoten
		dalam gelanggang $\End(X)$.
	\end{enumerate}
\end{lema}
% segment-id: o014.li2.chapter2.direct-sum-decomposition.pf005
\begin{bukti}
	Untuk $n\in\Z_{\geq 0}$, definisikan $X_n:=\Image(f^n)$, sehingga
	$X_0=X$. Dengan menerapkan Proposisi~\ref{prop:Images-f-gf} dan
	$\Coim\rightiso\Image$, diperoleh epimorfisme
	$\alpha_n:X_n\to X_{n+1}$ (yang diinduksi oleh $f$) dan monomorfisme
	$\beta_n:X_{n+1}\to X_n$. Dengan demikian, diperoleh rantai ganda
	$(X_n,\alpha_n,\beta_n)_{n=0}^\infty$.

	Untuk $n\gg 0$, $\alpha_n$ dan $\beta_n$ merupakan isomorfisme. Oleh karena
	itu, subobjek $\Image(f^\infty):=\Image(f^n)$ dari $X$, dengan $n\gg 0$,
	terdefinisi dengan baik; monomorfisme yang bersesuaian dilambangkan dengan
	$\iota:\Image(f^\infty)\hookrightarrow X$. Selain itu,
	$\Ker(f^n)=\Ker[X\to\Image(f^n)]$ juga merupakan subobjek
	$\Ker(f^\infty)$ yang tidak bergantung pada $n$ untuk $n\gg 0$.

	Dengan demikian, untuk $n\gg 0$, morfisme
	$\alpha_{2n-1}\cdots\alpha_n:X_n\to X_{2n}$ bersifat invertibel. Notasikan
	inversnya dengan
	$\psi:\Image(f^\infty)\rightiso\Image(f^\infty)$, dan tetapkan
	% segment-id: o014.li2.chapter2.direct-sum-decomposition.q008
	\[ p := \psi f^n: X \to \Image(f^\infty), \quad \Ker(p) = \Ker(f^\infty) \qquad (n \gg 0). \]
	Dari $p\iota=\identity_{\Image(f^\infty)}$ diperoleh bahwa
	$e:=\iota p\in\End(X)$ merupakan idempoten, dengan
	$\Image(e)=\Image(f^\infty)$ dan $\Ker(e)=\Ker(f^\infty)$. Dengan
	menerapkan Proposisi~\ref{prop:idempotent-decomposition}, diperoleh
	dekomposisi dalam (i).

	Jika $X$ tak terurai, maka $\Image(f^\infty)=0$ dan
	$\Ker(f^\infty)=X$, sehingga $f$ nilpoten; atau
	$\Image(f^\infty)=X$ dan $\Ker(f^\infty)=0$, sehingga $f$ invertibel.
	Inilah pernyataan (ii).
\end{bukti}

% segment-id: o014.li2.chapter2.direct-sum-decomposition.p012
Ingat dari \cite[Definisi 6.12.3]{Li1}: jika $S$ merupakan gelanggang dan
$S\smallsetminus S^\times$ merupakan ideal dua sisi, maka $S$ disebut
\emph{gelanggang lokal}.
\index{gelanggang lokal (local ring)}

% segment-id: o014.li2.chapter2.direct-sum-decomposition.co002
\begin{korolari}\label{prop:indecomposable-local-ring}
	Misalkan objek tak nol $X$ dalam suatu kategori abelian memenuhi syarat
	rantai ganda. Maka $X$ tak terurai jika dan hanya jika $\End(X)$ merupakan
	gelanggang lokal.
\end{korolari}
% segment-id: o014.li2.chapter2.direct-sum-decomposition.pf006
\begin{bukti}
	Lema~\ref{prop:Abel-cat-Fitting} menunjukkan bahwa jika $X$ tak terurai,
	setiap unsur $\End(X)$ bersifat nilpoten atau invertibel. Argumen selebihnya
	hanya menyangkut struktur gelanggang $\End(X)$ dan sama dengan
	\cite[Lema 6.12.6]{Li1}.
\end{bukti}

% segment-id: o014.li2.chapter2.direct-sum-decomposition.p013
Sekarang kita dapat menyatakan versi Teorema Krull--Remak--Schmidt untuk
kategori abelian.
% segment-id: o014.li2.chapter2.direct-sum-decomposition.th001
\begin{teorema}[M.\ Atiyah]\label{prop:Krull-Schmidt-gen}
	\index{Teorema Krull--Remak--Schmidt (Krull--Remak--Schmidt Theorem)}
	Misalkan $X$ merupakan objek tak nol dalam suatu kategori abelian.
	% segment-id: o014.li2.chapter2.direct-sum-decomposition.l005
	\begin{enumerate}[(i)]
		% segment-id: o014.li2.chapter2.direct-sum-decomposition.i014
		\item Jika $X$ memenuhi syarat rantai ganda, maka terdapat
		$n\in\Z_{\geq 1}$ dan subobjek-subobjek tak terurai
		$X_1,\ldots,X_n$ sedemikian sehingga
		$X=\bigoplus_{i=1}^n X_i$, dan setiap $\End(X_i)$ merupakan gelanggang
		lokal.
		% segment-id: o014.li2.chapter2.direct-sum-decomposition.i015
		\item Andaikan $X$ mempunyai dekomposisi $\bigoplus_{i=1}^n X_i$ dan
		$\bigoplus_{j=1}^m X'_j$, dengan setiap $X_i$ dan $X'_j$ tak terurai
		serta setiap $\End(X_i)$ dan $\End(X'_j)$ merupakan gelanggang lokal.
		Maka $n=m$, dan terdapat bijeksi $\sigma$ dari $\{1,\ldots,n\}$ ke
		dirinya sendiri sedemikian sehingga $X_i\simeq X'_{\sigma(i)}$ untuk
		setiap $i$.
	\end{enumerate}
\end{teorema}
% segment-id: o014.li2.chapter2.direct-sum-decomposition.pf007
\begin{bukti}
	Yang terutama perlu dibuktikan ialah bahwa $X$ yang memenuhi syarat rantai
	ganda pasti mempunyai dekomposisi seperti dalam (i). Bagian selebihnya hanya
	menyangkut operasi aljabar dalam gelanggang-gelanggang $\End$ dan sama dengan
	\cite[Teorema 6.12.8]{Li1}. Karena itu, untuk bagian berikut kita asumsikan
	bahwa syarat rantai ganda berlaku.

	Perhatikan bahwa jika $X=Y\oplus Z$ dengan $Y\neq 0$, maka $Y$ juga
	memenuhi syarat rantai ganda. Memang, untuk rantai ganda
	$(Y_n,\alpha_n,\beta_n)_{n=0}^\infty$ dengan $Y_0=Y$, ambil
	$X_n:=Y_n\oplus Z$, $\tilde{\alpha}_n:=\alpha_n\oplus\identity_Z$, dan
	$\tilde{\beta}_n:=\beta_n\oplus\identity_Z$. Dengan demikian diperoleh
	rantai ganda dengan $X_0=X$, sedangkan
	Lema~\ref{prop:isomorphism-matrix} menunjukkan bahwa
	$\tilde{\alpha}_n$ (atau $\tilde{\beta}_n$) merupakan isomorfisme jika dan
	hanya jika $\alpha_n$ (atau $\beta_n$) juga demikian.

	Jika $X$ sudah tak terurai, Korolari~\ref{prop:indecomposable-local-ring}
	menunjukkan bahwa $\End(X)$ merupakan gelanggang lokal; inilah (i). Jika
	(i) tidak berlaku untuk $X$, maka pasti terdapat subobjek tak nol
	$X_1,Y_1$ sedemikian sehingga $X=X_1\oplus Y_1$ dan pernyataan (i) tidak
	berlaku untuk $X_1$. Mengulangi proses ini pada $X_1$ memberikan
	$X_1=X_2\oplus Y_2$; iterasi menghasilkan rantai ganda
	$(X_n,\alpha_n,\beta_n)_{n=0}^\infty$, dengan
	$\alpha_n:X_n\twoheadrightarrow X_{n+1}$ dan
	$\beta_n:X_{n+1}\hookrightarrow X_n$ berasal dari
	$X_n=X_{n+1}\oplus Y_{n+1}$. Tidak satu pun morfisme ini merupakan
	isomorfisme dan $X_0:=X$, bertentangan dengan syarat rantai ganda. Ini
	membuktikan pernyataan tersebut.
\end{bukti}

% segment-id: o014.li2.chapter2.direct-sum-decomposition.p014
Dalam penerapan, yang penting adalah mengetahui kapan syarat rantai ganda
berlaku. Berikut ini salah satu syarat yang mencukupi.

% segment-id: o014.li2.chapter2.direct-sum-decomposition.pr003
\begin{proposisi}
	Misalkan $X$ merupakan objek berpanjang hingga dalam suatu kategori abelian
	(Definisi~\ref{def:finite-length-object}). Maka $X$ memenuhi syarat rantai
	ganda.
\end{proposisi}
% segment-id: o014.li2.chapter2.direct-sum-decomposition.pf008
\begin{bukti}
	Ambil rantai ganda $(X_n,\alpha_n,\beta_n)_{n=0}^\infty$ dengan
	$X_0=X$. Maka
	\[
		\left(\Image(\beta_0\cdots\beta_n)\right)_{n=0}^\infty
	\]
	merupakan rantai
	menurun dalam himpunan terurut parsial $\mathrm{Sub}_X$. Karena setiap
	$\beta_n$ merupakan monomorfisme, untuk $n\gg 0$ syarat Artinian menjamin
	adanya diagram komutatif
	% segment-id: o014.li2.chapter2.direct-sum-decomposition.g004
	\[\begin{tikzcd}
		X_{n+1} \arrow[d, "\beta_n"'] \arrow[r, "\sim"] & \Image(\beta_0 \cdots \beta_n) \arrow[d, "\sim" sloped] \arrow[hookrightarrow, r] & X_0 \\
		X_n \arrow[r, "\sim"] & \Image(\beta_0 \cdots \beta_{n-1}) \arrow[hookrightarrow, ru] &
	\end{tikzcd}\]
	dan $\beta_n$ merupakan isomorfisme. Demikian pula, terapkan syarat
	Noetherian pada rantai menaik
	$\left(\Ker(\alpha_n\cdots\alpha_0)\right)_{n=0}^\infty$ dan gunakan
	bahwa setiap $\alpha_n$ merupakan epimorfisme. Untuk $n\gg 0$, diperoleh
	diagram komutatif
	dengan baris-baris eksak
	% segment-id: o014.li2.chapter2.direct-sum-decomposition.g005
	\[\begin{tikzcd}
		0 \arrow[r] & \Ker(\alpha_{n-1} \cdots \alpha_0) \arrow[d, "\sim" sloped] \arrow[r] & X_0 \arrow[r, "{\alpha_{n-1} \cdots \alpha_0}" inner sep=0.7em] \arrow[equal, d] & X_n \arrow[d, "{\alpha_n}"] \arrow[r] & 0 \\
		0 \arrow[r] & \Ker(\alpha_n \cdots \alpha_0) \arrow[r] & X_0 \arrow[r, "{\alpha_n \cdots \alpha_0}"' inner sep=0.7em] & X_{n+1} \arrow[r] & 0
	\end{tikzcd}\]
	dan $\alpha_n$ merupakan isomorfisme.
\end{bukti}

% segment-id: o014.li2.chapter2.direct-sum-decomposition.p015
Bagian latihan akan memperkenalkan cara-cara lain untuk memperoleh syarat
rantai ganda.
